% =============================================================================
% PREÂMBULO INTEGRAL CUSTOMIZADO
% =============================================================================

% --- Pacotes de Geometria e Tipografia ---
\usepackage[a4paper,lmargin=1.5cm,rmargin=1.5cm,tmargin=1.2cm,bmargin=1.5cm]{geometry}
\usepackage[brazilian]{babel}
\usepackage[T1]{fontenc}
\usepackage{microtype}
\usepackage{xspace}

% --- Macros de Tipografia (Corrigindo T1) ---
\newcommand{\f}{\textordfeminine\xspace} % Para usar ª
\newcommand{\m}{\textordmasculine\xspace} % Para usar º

% --- Matemática e Símbolos ---
\usepackage{amsmath,amsfonts,amssymb,mathtools,mathrsfs,amsthm}
\usepackage{stmaryrd}
\usepackage{latexsym}
\usepackage{esint}

% --- Tikz e Referências ---
\usepackage{csquotes} % Elimina o aviso do babel
\usepackage{caption}
\usepackage{tikz}
\usetikzlibrary{arrows.meta, positioning, shapes.geometric, decorations.pathreplacing}
\usepackage{pgfplots}
\pgfplotsset{compat=1.18}
\tikzset{
	passo/.style={
		rectangle, 
		draw, 
		align=center, % Alinhamento centralizado para evitar hifenização forçada
		text width=6cm, % Largura otimizada para caber os nomes sem estourar
		rounded corners, 
		minimum height=1.2cm, 
		font=\small
	},
	arrow/.style={-Stealth, thick}
}
% Referências Bibliográficas
\usepackage[backend=biber, style=abnt]{biblatex} % Recomendado usar biber para bibliografias modernas
%\addbibresource{referencias.bib}

% --- Elementos Visuais e Layout ---
\usepackage[most]{tcolorbox}
\usepackage{enumitem}
\setlist[itemize]{noitemsep, topsep=0pt, parsep=0pt, partopsep=0pt}
\usepackage[normalem]{ulem}
\usepackage[Smaller]{cancel}
\usepackage{xcolor}
\usepackage[nodisplayskipstretch]{setspace}
\usepackage{etoolbox} % ESSENCIAL para os interruptores de blocos

% --- Navegação e Links ---
\usepackage[colorlinks=true, linkcolor=black, urlcolor=blue, hypertexnames=false]{hyperref}
\usepackage{bookmark}

% --- Configurações Globais de Estilo ---
\setstretch{1.3}
\setlength{\parindent}{0pt}
\relpenalty=10000
\binoppenalty=10000
\hbadness=10000
\everymath{\displaystyle}
\pagestyle{empty}
\newcommand{\linhaestilizada}{ % Linha centralizada verticalmente com diamond
	\par\noindent\makebox[\textwidth]{\rule[0.5ex]{0.47\textwidth}{0.4pt}\hfill$\diamond$\hfill\rule[0.5ex]{0.47\textwidth}{0.4pt}}\par
}
\setlength{\overfullrule}{5pt} % Desenha uma barra de 5pt no final de hboxes para erros de estrapolação de linha horizontal

% =============================================================================
% SISTEMA DE CONTROLE DE VISIBILIDADE (INTERRUPTORES)
% =============================================================================
% Altere 'true' para 'false' para omitir o bloco no PDF final sem apagar o código.

\newbool{showanalise}
\setbool{showanalise}{true}

\newbool{showsintese}
\setbool{showsintese}{true}

\newbool{showaprofundamento} 
\setbool{showaprofundamento}{true}

\newbool{showtheorem}
\setbool{showtheorem}{true}

\newbool{showproof}
\setbool{showproof}{true}

\newbool{showtextoauxiliar}
\setbool{showtextoauxiliar}{true} 

% =============================================================================
% COMANDOS ESTRUTURAIS DO PROJETO
% =============================================================================

% Valor padrão seguro para compilação via pacote subfiles
\providecommand{\listid}{L0}

% Define o ambiente theorem* como não numerado
% --- Ambientes Condicionais de Teoremas e Provas ---
\makeatletter
\@ifundefined{theorem*}{%
	\theoremstyle{definition} 
	\newtheorem*{theorem*}{Teorema}%
}{}
\makeatother

\newenvironment{teorema}{%
	\ifbool{showtheorem}{%
		\begin{theorem*}%
		}{%
			\setbox0=\vbox\bgroup
		}%
	}{%
		\ifbool{showtheorem}{%
		\end{theorem*}%
	}{%
		\egroup
	}%
}

\newenvironment{prova}{%
	\ifbool{showproof}{%
		\begin{proof}[\normalfont\underline{Demonstração}:] 
		}{%
			\setbox0=\vbox\bgroup
		}%
	}{%
		\ifbool{showproof}{%
		\end{proof}%
	}{%
		\egroup
	}%
}

% --- Comando Customizado: Texto Auxiliar Solto Condicional ---
\newcommand{\textoauxiliar}[1]{%
	\ifbool{showtextoauxiliar}{%
		\par\vspace{3pt}%
		#1% Renderiza apenas o texto solto, sem cabeçalho
		\par\vspace{3pt}%
	}{}%
}

\newcounter{numquestao}

% =============================================================================
% SISTEMA INTELIGENTE DE SELEÇÃO DE QUESTÕES
% =============================================================================

% 1. Controle de Estados
\newbool{showcurrentQ}
\setbool{showcurrentQ}{true} 

\newbool{modoisolado}
\setbool{modoisolado}{false} 

% 2. Comandos de Seleção Otimizados para Listas (Separados por vírgula)
\newcommand{\ocultarquestao}[1]{%
	\foreach \q in {#1} {%
		\csgdef{hideQ\q}{true}% Usa csgdef global para escapar o escopo do foreach do tikz
	}%
}

\newcommand{\mostrarapenasquestao}[1]{%
	\setbool{modoisolado}{true}%
	\foreach \q in {#1} {%
		\csgdef{showOnlyQ\q}{true}%
	}%
}

% 3. Sobrescrita do Comando Questão
\newcommand{\questao}[1]{%
	\stepcounter{numquestao}% Avança o contador SEMPRE, mantendo a numeração original 
	%
	% Lógica de Visibilidade Dinâmica
	\setbool{showcurrentQ}{true}%
	\ifbool{modoisolado}{%
		\ifcsdef{showOnlyQ\arabic{numquestao}}{}{\setbool{showcurrentQ}{false}}%
	}{%
		\ifcsdef{hideQ\arabic{numquestao}}{\setbool{showcurrentQ}{false}}{}%
	}%
	%
	% Renderização Condicional
	\ifbool{showcurrentQ}{%
		\hypertarget{dest.\listid.\arabic{numquestao}}{}%
		\bookmark[level=1, dest={dest.\listid.\arabic{numquestao}}]{Questão \arabic{numquestao}}%
		\begin{tcolorbox}[
			enhanced, sharp corners, boxrule=0.5pt,
			colback=gray!6!white, colframe=black,         
			left=1pt, right=1pt, top=1pt, bottom=1pt,
			width=1\textwidth, before skip=3pt, after skip=3pt,
			fontupper=\normalfont
			]%
			\mbox{\textbf{Questão \arabic{numquestao}.}}\hspace{0.3em}\ignorespaces#1%
		\end{tcolorbox}%
	}{}%
}

% 4. Condicionais em cascata para os blocos da resolução 
\newcommand{\analise}[1]{%
	\ifbool{showcurrentQ}{%
		\ifbool{showanalise}{%
			\par\vspace{5pt}\textbf{\underline{Análise e Estratégia}}:\par 
			#1\par
		}{}%
	}{}%
}

\newcommand{\solucao}[1]{%
	\ifbool{showcurrentQ}{%
		\par\vspace{5pt}
		\textbf{\underline{Solução}}:\par
		#1 \hfill$\blacksquare$
		\par
	}{}%
}

\newcommand{\solucaoalt}[1]{%
	\ifbool{showcurrentQ}{%
		\ifbool{showanalise}{%
			\par\vspace{5pt}
			\textbf{\underline{Solução Alternativa}}:\par 
			#1 \hfill$\blacksquare$
			\par
		}{}%
	}{}%
}

\newcommand{\sintese}[1]{%
	\ifbool{showcurrentQ}{%
		\ifbool{showsintese}{%
			\par\vspace{5pt}\textbf{\underline{Síntese da Construção}}:\par 
			#1\par
		}{}%
	}{}%
}

\newcommand{\aprofundamento}[1]{%
	\ifbool{showcurrentQ}{%
		\ifbool{showaprofundamento}{%
			\linhaestilizada
			\par\noindent\textbf{\underline{Aprofundamento Teórico}}\par
			#1
			\par
			\linhaestilizada
		}{}%
	}{}%
}

% =============================================================================
% OPERADORES E NOTAÇÃO DE ANÁLISE
% =============================================================================
\newcommand{\R}{\mathbb{R}}
\newcommand{\norm}[1]{\|#1\|}
\newcommand{\inner}[2]{\langle #1, #2 \rangle}
\newcommand{\Sobolev}[2]{W^{#1, #2}}
\newcommand{\Hspace}[1]{H^#1}
\DeclareMathOperator{\supp}{supp}
\DeclareMathOperator{\dist}{dist}
\newcommand{\FF}{\mathcal{F}} 
\newcommand{\invFF}{\mathcal{F}^{-1}}